\documentclass[../../main.tex]{subfiles}

\begin{document}

\section{$L^{p}$ 是完备的: Riesz-Fischer 定理}\label{section:ch07-Riesz-Fischer}

采用与赋予绝对值范数的 $\mathbf{R}$ 中的序列完全相同的方式, 对赋范线性空间中的序列定义收敛序列与 Cauchy 序列.

\begin{definition}[收敛与 Cauchy 序列]\label{definition:ch07-收敛与Cauchy}
赋予范数 $\|\cdot\|$ 的线性空间 $X$ 中的序列 $\{f_{n}\}$ 称为在 $X$ 中收敛于 $f$, 若
\begin{align*}
\lim_{n\to\infty}\|f-f_{n}\|=0.
\end{align*}
我们记作:
\begin{align*}
\text{在 }X\text{ 中 }\{f_{n}\}\to f,\quad\text{或在 }X\text{ 中 }\lim_{n\to\infty}f_{n}=f,
\end{align*}
意味着每个 $f_{n}$ 和 $f$ 属于 $X$ 且 $\lim_{n\to\infty}\|f-f_{n}\|=0$.
\end{definition}

显然对于 $C[a,b]$ 中的序列 $\{f_{n}\}$ 和函数 $f$, 在赋予最大值范数的 $C[a,b]$ 中, $\{f_{n}\}\to f$ 当且仅当在 $[a,b]$ 上一致地 $\{f_{n}\}\to f$. 此外, 由于 $L^{\infty}(E)$ 中的函数的本性上确界是一个本性上界, 对于 $L^{\infty}(E)$ 中的序列 $\{f_{n}\}$ 和函数 $f$, 在 $L^{\infty}(E)$ 中 $\{f_{n}\}\to f$ 当且仅当在零测度集的补集上一致地 $\{f_{n}\}\to f$. 对于 $L^{p}(E)$ 中的序列 $\{f_{n}\}$ 和函数 $f$, $1\leqslant p<\infty$, 在 $L^{p}(E)$ 中 $\{f_{n}\}\to f$ 当且仅当
\begin{align*}
\lim_{n\to\infty}\int_{E}|f_{n}-f|^{p}=0.
\end{align*}

\begin{definition}[Cauchy 序列与 Banach 空间]\label{definition:ch07-Cauchy与Banach}
赋予范数 $\|\cdot\|$ 的线性空间 $X$ 中的序列 $\{f_{n}\}$ 称为 $X$ 中的 Cauchy 序列, 若对每个 $\varepsilon>0$, 存在自然数 $N$ 使得对所有 $m$, $n\geqslant N$, $\|f_{n}-f_{m}\|<\varepsilon$.

赋范线性空间 $X$ 称为完备的, 若 $X$ 中的每个 Cauchy 序列收敛到 $X$ 中的某个函数. 完备的赋范线性空间称为 Banach 空间.
\end{definition}

实数的完备性公理等价于赋予绝对值范数的 $\mathbf{R}$ 是完备的这一断言. 这立即蕴涵欧氏空间 $\mathbf{R}^{n}$ 也是完备的. 在数学分析的基础课程中证明了赋予最大值范数的 $C[a,b]$ 是完备的 (见习题 31). 同样的方法与可测函数的逐点极限的可测性一起, 表明 $L^{\infty}(E)$ 也是完备的 (见习题 33).

\begin{proposition}\label{proposition:ch07-收敛序列Cauchy}
令 $X$ 为赋范线性空间. 则 $X$ 中的每个收敛的序列是 Cauchy 的. 此外, $X$ 中的 Cauchy 序列收敛, 若它有一个收敛的子序列.
\end{proposition}

\begin{proof}
令在 $X$ 中 $\{f_{n}\}\to f$. 根据范数的三角不等式, 对所有 $m$, $n$,
\begin{align*}
\|f_{n}-f_{m}\|=\|[f_{n}-f]+[f-f_{m}]\|\leqslant\|f_{n}-f\|+\|f_{m}-f\|.
\end{align*}
因此 $\{f_{n}\}$ 是 Cauchy 的.

现在令 $\{f_{n}\}$ 为 $X$ 中的 Cauchy 序列. 它有在 $X$ 中收敛到 $f$ 的子序列 $\{f_{n_{k}}\}$. 令 $\varepsilon>0$. 由于 $\{f_{n}\}$ 是 Cauchy 的, 我们可以选取 $N$ 使得对所有 $m$, $n\geqslant N$, $\|f_{n}-f_{m}\|_{p}<\varepsilon/2$. 由于 $\{f_{n_{k}}\}$ 收敛到 $f$, 可以选取 $k$ 使得 $n_{k}>N$ 且 $\|f_{n_{k}}-f\|<\varepsilon/2$. 则根据范数的三角不等式, 对 $n\geqslant N$,
\begin{align*}
\|f_{n}-f\|=\|[f_{n}-f_{n_{k}}]+[f_{n_{k}}-f]\|=\|f_{n}-f_{n_{k}}\|+\|f_{n_{k}}-f\|<\varepsilon.
\end{align*}
因此在 $X$ 中 $\{f_{n}\}\to f$.
\end{proof}

根据以上命题, 证明特定赋范线性空间的完备性的有用策略是证明针对空间的性质, 特定类型的 Cauchy 序列收敛, 并且证明每个 Cauchy 序列有这种特定类型的子序列 (见习题 30 和 32). 对 $1\leqslant p<\infty$, 在 $L^{p}(E)$ 空间中, 定义如下的快速 Cauchy 序列\footnote{在论文 ``Rethinking the Lebesgue Integral'' (美国数学月刊, 2009 年 12 月), Peter Lax 单独挑出连续函数序列的逐点极限作为构造完备空间 $L^{1}$ 的主要对象, 该连续函数序列是关于 $L^{1}$ 范数的快速 Cauchy 序列. 他捍卫了以下观点: 在积分的研究中使用关于 Banach 空间的定理, 其主要目标是对 $L^{1}$ 的识别. 无须先对测度论做单独的研究, Lax 构造了 $L^{1}$ 中的函数作为连续函数的快速 Cauchy 序列的极限.}是有用的.

\begin{definition}[快速 Cauchy 序列]\label{definition:ch07-快速Cauchy}
令 $X$ 为具有范数 $\|\cdot\|$ 的赋范线性空间. $X$ 中的序列 $\{f_{n}\}$ 称为快速 Cauchy 序列, 若存在收敛的正项级数 $\sum_{k=1}^{\infty}\varepsilon_{k}$, 使得
\begin{align*}
\text{对所有 }k,\quad\|f_{k+1}-f_{k}\|\leqslant\varepsilon_{k}^{2}.
\end{align*}
\end{definition}

以下观察是有用的: 若 $\{f_{n}\}$ 是赋范线性空间中的序列, 而非负数序列 $\{a_{k}\}$ 具有对所有 $k$,
\begin{align*}
\|f_{k+1}-f_{k}\|\leqslant a_{k}
\end{align*}
的性质, 则由于
\begin{align*}
\text{对所有 }n,k,\quad f_{n+k}-f_{n}=\sum_{j=n}^{n+k-1}[f_{j+1}-f_{j}],
\end{align*}
\begin{align}
\text{对所有 }n,k,\quad\|f_{n+k}-f_{n}\|\leqslant\sum_{j=n}^{n+k-1}\|f_{j+1}-f_{j}\|\leqslant\sum_{j=n}^{\infty}a_{j}.\label{align:ch07-快速Cauchy估计}
\end{align}

\begin{proposition}\label{proposition:ch07-快速Cauchy}
令 $X$ 为赋范线性空间. 则 $X$ 中的每个快速 Cauchy 序列是 Cauchy 的. 此外, 每个 Cauchy 序列有一个快速 Cauchy 的子序列.
\end{proposition}

\begin{proof}
令 $\{f_{n}\}$ 为 $X$ 中的快速 Cauchy 序列, 而 $\sum_{k=1}^{\infty}\varepsilon_{k}$ 是收敛的正项级数, 使得对所有 $k$,
\begin{align}
\|f_{k+1}-f_{k}\|\leqslant\varepsilon_{k}^{2}.\label{align:ch07-快速Cauchy条件}
\end{align}
我们从\eqref{align:ch07-快速Cauchy估计}推出对所有 $n$, $k$,
\begin{align}
\|f_{n+k}-f_{n}\|\leqslant\sum_{j=n}^{\infty}\varepsilon_{j}^{2}.\label{align:ch07-快速Cauchy推论}
\end{align}
由于级数 $\sum_{k=1}^{\infty}\varepsilon_{k}$ 收敛, 级数 $\sum_{k=1}^{\infty}\varepsilon_{k}^{2}$ 也收敛. 我们从\eqref{align:ch07-快速Cauchy推论}推出 $\{f_{n}\}$ 是 Cauchy 的. 现在假设 $\{f_{n}\}$ 是 $X$ 中的 Cauchy 序列. 我们可以归纳地选取严格递增的自然数列 $\{n_{k}\}$, 使得对所有 $k$,
\begin{align*}
\|f_{n_{k+1}}-f_{n_{k}}\|\leqslant(1/2)^{k}.
\end{align*}
因为公比为 $1/\sqrt{2}$ 的几何级数收敛, 所以子序列 $\{f_{n_{k}}\}$ 是快速 Cauchy 的.
\end{proof}

\begin{theorem}\label{theorem:ch07-快速Cauchy逐点收敛}
令 $E$ 为可测集而 $1\leqslant p\leqslant\infty$. 则 $L^{p}(E)$ 中的每个快速 Cauchy 序列关于 $L^{p}(E)$ 范数在 $E$ 上 a.e. 逐点地收敛到 $L^{p}(E)$ 中的函数 $f$.
\end{theorem}

\begin{proof}
我们将 $p=\infty$ 的情形留作练习 (习题 33). 假设 $1\leqslant p<\infty$. 令 $\{f_{n}\}$ 为 $L^{p}(E)$ 中快速收敛的序列. 通过从 $E$ 中去掉一个测度为零的集合, 我们可以假设每个 $f_{n}$ 取实值. 选取收敛的正项级数 $\sum_{k=1}^{\infty}\varepsilon_{k}$, 使得对所有 $k$,
\begin{align}
\|f_{k+1}-f_{k}\|_{p}\leqslant\varepsilon_{k}^{2}.\label{align:ch07-快速Cauchy范数}
\end{align}
因此对所有 $k$,
\begin{align}
\int_{E}|f_{k+1}-f_{k}|^{p}\leqslant\varepsilon_{k}^{2p}.\label{align:ch07-快速Cauchy积分}
\end{align}
固定一个自然数 $k$. 由于对 $x\in E$, $|f_{k+1}(x)-f_{k}(x)|\geqslant\varepsilon_{k}$ 当且仅当 $|f_{k+1}(x)-f_{k}(x)|^{p}\geqslant\varepsilon_{k}^{p}$, 我们从\eqref{align:ch07-快速Cauchy积分}和 Chebyshev 不等式推出,
\begin{align*}
m\{x\in E\mid|f_{k+1}(x)-f_{k}(x)|\geqslant\varepsilon_{k}\}=m\{x\in E\mid|f_{k+1}(x)-f_{k}(x)|^{p}\geqslant\varepsilon_{k}^{p}\}\leqslant\frac{1}{\varepsilon_{k}^{p}}\cdot\int_{E}|f_{k+1}-f_{k}|^{p}\leqslant\varepsilon_{k}^{p}.
\end{align*}
由于 $p\geqslant 1$, 级数 $\sum_{k=1}^{\infty}\varepsilon_{k}^{p}$ 收敛. Borel-Cantelli 引理告诉我们存在 $E$ 的子集 $E_{0}$, 它的测度为零, 且对每个 $x\in E\sim E_{0}$ 存在指标 $K(x)$, 使得对所有 $k\geqslant K(x)$,
\begin{align*}
|f_{k+1}(x)-f_{k}(x)|<\varepsilon_{k}.
\end{align*}
令 $x$ 属于 $E\sim E_{0}$. 则对所有 $n\geqslant K(x)$ 和所有 $k$,
\begin{align}
|f_{n+k}(x)-f_{n}(x)|\leqslant\sum_{j=n}^{n+k-1}|f_{j+1}(x)-f_{j}(x)|\leqslant\sum_{j=n}^{\infty}\varepsilon_{j}.\label{align:ch07-逐点Cauchy}
\end{align}
级数 $\sum_{j=1}^{\infty}\varepsilon_{j}$ 收敛, 因此实数序列 $\{f_{k}(x)\}$ 是 Cauchy 的. 实数是完备的. 记 $\{f_{k}(x)\}$ 的极限为 $f(x)$. 在 $E_{0}$ 上定义 $f=0$. 从\eqref{align:ch07-快速Cauchy范数}和\eqref{align:ch07-快速Cauchy估计}推出对所有 $n$, $k$,
\begin{align}
\int_{E}|f_{n+k}-f_{n}|^{p}\leqslant\left[\sum_{j=n}^{\infty}\varepsilon_{j}^{2}\right]^{p}.\label{align:ch07-范数估计}
\end{align}
由于在 $E$ 上 a.e. 逐点 $\{f_{n}\}\to f$, 在该不等式中取 $k\to\infty$ 时的极限, 且从 Fatou 引理推出, 对所有 $n$,
\begin{align*}
\int_{E}|f-f_{n}|^{p}\leqslant\left[\sum_{j=n}^{\infty}\varepsilon_{j}^{2}\right]^{p}.
\end{align*}
由于级数 $\sum_{k=1}^{\infty}\varepsilon_{k}^{2}$ 收敛, $f$ 属于 $L^{p}(E)$ 且在 $L^{p}(E)$ 中 $\{f_{n}\}\to f$. 我们构造了 $f$ 作为 $\{f_{n}\}$ 在 $E$ 上的 a.e. 逐点极限, 证明完毕.
\end{proof}

\begin{theorem}[Riesz-Fischer 定理]\label{theorem:ch07-Riesz-Fischer}
令 $E$ 为可测集而 $1\leqslant p\leqslant\infty$. 则 $L^{p}(E)$ 是 Banach 空间. 此外, 若在 $L^{p}(E)$ 中 $\{f_{n}\}\to f$, 则 $\{f_{n}\}$ 的一个子序列在 $E$ 上 a.e. 逐点收敛到 $f$.
\end{theorem}

\begin{proof}
令 $\{f_{n}\}$ 为 $L^{p}(E)$ 中的 Cauchy 序列. 根据命题 5, 存在 $\{f_{n}\}$ 的子序列 $\{f_{n_{k}}\}$, 它是快速 Cauchy 的. 前一个定理告诉我们 $\{f_{n_{k}}\}$ 关于 $L^{p}(E)$ 范数在 $E$ 上 a.e. 逐点收敛到 $L^{p}(E)$ 中的函数 $f$. 根据命题 4 整个 Cauchy 序列关于 $L^{p}(E)$ 范数收敛到 $f$. 最后, 若序列在 $L^{p}(E)$ 中收敛, 则它有快速 Cauchy 子序列且该子序列在 $E$ 上 a.e. 逐点收敛到同样的极限函数.
\end{proof}

如同以下例子所表明的, $L^{p}(E)$ 中的序列 $\{f_{n}\}$ 在 $E$ 上 a.e. 逐点收敛到 $L^{p}(E)$ 中的 $f$, 一般不在 $L^{p}(E)$ 中收敛.

\begin{example}\label{example:ch07-逐点不收敛}
对于 $E=[0,1]$, $1\leqslant p<\infty$ 和每个自然数 $n$, 令 $f_{n}=n^{1/p}\chi_{(0,1/n)}$. 该序列在 $[0,1]$ 上逐点收敛到恒等于零的函数, 但不关于 $L^{p}[0,1]$ 范数收敛到该函数.
\end{example}

以下两个定理提供了逐点收敛的序列在 $L^{p}(E)$ 中收敛的必要与充分条件.

\begin{theorem}\label{theorem:ch07-范数收敛等价}
令 $E$ 为可测集而 $1\leqslant p<\infty$. 假定 $\{f_{n}\}$ 是 $L^{p}(E)$ 中在 $E$ 上 a.e. 逐点收敛于属于 $L^{p}(E)$ 的函数 $f$ 的序列. 则
\begin{align*}
\text{在 }L^{p}(E)\text{ 中 },\{f_{n}\}\to f\text{ 当且仅当 }\lim_{n\to\infty}\int_{E}|f_{n}|^{p}=\int_{E}|f|^{p}.
\end{align*}
\end{theorem}

\begin{proof}
通过可能地从 $E$ 中去掉一个测度为零的集合, 我们可以假设 $f$ 和每个 $f_{n}$ 在整个 $E$ 上是实值的且收敛是逐点的. 我们从 Minkowski 不等式推出, 对每个 $n$, $|\|f_{n}\|_{p}-\|f\|_{p}|\leqslant\|f_{n}-f\|_{p}$. 因此, 若在 $L^{p}(E)$ 中 $\{f_{n}\}\to f$, 则 $\lim_{n\to\infty}\int_{E}|f_{n}|^{p}=\int_{E}|f|^{p}$. 为证明反命题, 假设 $\lim_{n\to\infty}\int_{E}|f_{n}|^{p}=\int_{E}|f|^{p}$. 对所有 $t$ 定义 $\psi(t)=t^{p}$. 则 $\psi$ 是凸的, 这是由于它的二阶导数是非负的, 且因此对所有 $a$, $b$,
\begin{align*}
\psi\left(\frac{a+b}{2}\right)\leqslant\frac{\psi(a)+\psi(b)}{2}.
\end{align*}
因此,
\begin{align*}
\text{对所有 }a,b,0\leqslant\frac{|a|^{p}+|b|^{p}}{2}-\left|\frac{a-b}{2}\right|^{p}.
\end{align*}
因此, 对每个 $n$, $E$ 上的非负可测函数 $h_{n}$ 定义为, 对所有 $x\in E$,
\begin{align*}
h_{n}(x)=\frac{|f_{n}(x)|^{p}+|f(x)|^{p}}{2}-\left|\frac{f_{n}(x)-f(x)}{2}\right|^{p}.
\end{align*}
由于在 $E$ 上逐点地 $\{h_{n}\}\to|f|^{p}$, 我们从 Fatou 引理推出
\begin{align*}
\int_{E}|f|^{p}&\leqslant\liminf\left[\int_{E}h_{n}\right]\\
&=\liminf\left[\int_{E}\frac{|f_{n}(x)|^{p}+|f(x)|^{p}}{2}-\left|\frac{f_{n}(x)-f(x)}{2}\right|^{p}\right]\\
&=\int_{E}|f|^{p}-\limsup\left[\int_{E}\left|\frac{f_{n}(x)-f(x)}{2}\right|^{p}\right].
\end{align*}
因此,
\begin{align*}
\limsup\left[\int_{E}\left|\frac{f_{n}(x)-f(x)}{2}\right|^{p}\right]\leqslant 0,
\end{align*}
即在 $L^{p}(E)$ 中 $\{f_{n}\}\to f$.
\end{proof}

\begin{theorem}\label{theorem:ch07-Vitali判别}
令 $E$ 为可测集而 $1\leqslant p<\infty$. 假定 $\{f_{n}\}$ 是 $L^{p}(E)$ 中在 $E$ 上 a.e. 逐点收敛于属于 $L^{p}(E)$ 的函数 $f$ 的序列. 则
\begin{align*}
\text{在 }L^{p}(E)\text{ 中 },\{f_{n}\}\to f\text{ 当且仅当 }\{|f_{n}|^{p}\}\text{ 在 }E\text{ 上是一致可积与紧的}.
\end{align*}
\end{theorem}

\begin{proof}
非负可积序列 $\{|f_{n}-f|^{p}\}$ 在 $E$ 上 a.e. 逐点收敛到零. 根据第 5 章的推论 2——Vitali 收敛定理的一个推论,
\begin{align*}
\lim_{n\to\infty}\int_{E}|f_{n}-f|^{p}=0\text{ 当且仅当在 }E\text{ 上 }\{|f_{n}-f|^{p}\}\text{ 是一致可积与紧的}.
\end{align*}
然而, 我们从不等式\eqref{align:ch07-三角不等式估计}推出对所有 $n$, 在 $E$ 上 a.e.
\begin{align*}
|f_{n}-f|^{p}\leqslant 2^{p}\{|f_{n}|^{p}+|f|^{p}\}\text{ 且 }|f_{n}|^{p}\leqslant 2^{p}\{|f_{n}-f|^{p}+|f|^{p}\}.
\end{align*}
根据假设, $|f|^{p}$ 在 $E$ 上可积, 因此 $\{|f_{n}-f|^{p}\}$ 在 $E$ 上是一致可积与紧的当且仅当序列 $\{|f_{n}|^{p}\}$ 在 $E$ 上是一致可积与紧的.
\end{proof}

\begin{exercise}\label{exercise:ch07-23}
提供一个具有以下性质的实数序列的例子: 它是 Cauchy 的但不是快速 Cauchy 的.
\end{exercise}

\begin{exercise}\label{exercise:ch07-24}
令 $X$ 为赋范线性空间. 假设在 $X$ 中 $\{f_{n}\}\to f$, $\{g_{n}\}\to g$, 而 $\alpha$ 和 $\beta$ 是实数. 证明在 $X$ 中,
\begin{align*}
\{\alpha f_{n}+\beta g_{n}\}\to\alpha f+\beta g.
\end{align*}
\end{exercise}

\begin{exercise}\label{exercise:ch07-25}
假设 $E$ 有有限测度而 $1\leqslant p_{1}<p_{2}\leqslant\infty$. 证明: 若在 $L^{p_{2}}(E)$ 中 $\{f_{n}\}\to f$, 则在 $L^{p_{1}}(E)$ 中 $\{f_{n}\}\to f$.
\end{exercise}

\begin{exercise}\label{exercise:ch07-26}
($L^{p}$ 控制收敛定理) 令 $\{f_{n}\}$ 为在 $E$ 上 a.e. 逐点收敛于 $f$ 的可测函数序列. 对 $1\leqslant p<\infty$, 假定存在 $L^{p}(E)$ 中的函数 $g$, 使得对所有 $n$, 在 $E$ 上 a.e. $|f_{n}|\leqslant g$. 证明在 $L^{p}(E)$ 中 $\{f_{n}\}\to f$.
\end{exercise}

\begin{exercise}\label{exercise:ch07-27}
$E$ 为可测集而 $1\leqslant p<\infty$. 假设在 $L^{p}(E)$ 中 $\{f_{n}\}\to f$. 证明存在子序列 $\{f_{n_{k}}\}$ 和函数 $g\in L^{p}(E)$, 使得在 $E$ 上对所有 $k$, a.e. $|f_{n_{k}}|\leqslant g$.
\end{exercise}

\begin{exercise}\label{exercise:ch07-28}
假设 $E$ 有有限测度而 $1\leqslant p<\infty$. 假定 $\{f_{n}\}$ 是在 $E$ 上 a.e. 逐点收敛于 $f$ 的可测函数序列. 对 $1\leqslant p<\infty$, 证明: 若存在 $\theta>0$, 使得 $\{f_{n}\}$ 属于 $L^{p+\theta}(E)$ 且作为 $L^{p+\theta}(E)$ 的子集是有界的, 则在 $L^{p}(E)$ 中 $\{f_{n}\}\to f$.
\end{exercise}

\begin{exercise}\label{exercise:ch07-29}
考虑 $[a,b]$ 上的赋予范数 $\|\cdot\|_{\max}$ 的多项式线性空间. 该赋范线性空间是 Banach 空间吗?
\end{exercise}

\begin{exercise}\label{exercise:ch07-30}
令 $\{f_{n}\}$ 为 $C[a,b]$ 中的序列, 而 $\sum_{k=1}^{\infty}a_{k}$ 是收敛的正项级数, 使得对所有 $k$, $\|f_{k+1}-f_{k}\|_{\max}\leqslant a_{k}$. 证明对所有 $k$, $n$ 和所有 $x\in[a,b]$,
\begin{align*}
|f_{n+k}(x)-f_{k}(x)|\leqslant\|f_{n+k}-f_{k}\|_{\max}\leqslant\sum_{j=n}^{\infty}a_{j}.
\end{align*}
得出存在函数 $f\in C[a,b]$, 使得在 $[a,b]$ 上一致地 $\{f_{n}\}\to f$.
\end{exercise}

\begin{exercise}\label{exercise:ch07-31}
用前一个习题证明赋予最大值范数的 $C[a,b]$ 是 Banach 空间.
\end{exercise}

\begin{exercise}\label{exercise:ch07-32}
令 $\{f_{n}\}$ 为 $L^{\infty}(E)$ 中的序列, 而 $\sum_{k=1}^{\infty}a_{k}$ 是收敛的正项级数, 使得对所有 $k$,
\begin{align*}
\|f_{k+1}-f_{k}\|_{\infty}\leqslant a_{k}.
\end{align*}
证明存在 $E$ 的子集 $E_{0}$, 它的测度为零且对所有 $k$, $n$ 和所有 $x\in E\sim E_{0}$,
\begin{align*}
|f_{n+k}(x)-f_{k}(x)|\leqslant\|f_{n+k}-f_{k}\|_{\infty}\leqslant\sum_{j=n}^{\infty}a_{j}.
\end{align*}
证明存在函数 $f\in L^{\infty}(E)$ 使得在 $E\sim E_{0}$ 上一致地 $\{f_{n}\}\to f$.
\end{exercise}

\begin{exercise}\label{exercise:ch07-33}
用前一个习题证明 $L^{\infty}(E)$ 是一个 Banach 空间.
\end{exercise}

\begin{exercise}\label{exercise:ch07-34}
证明对于 $1\leqslant p\leqslant\infty$, $\ell^{p}$ 是 Banach 空间.
\end{exercise}

\begin{exercise}\label{exercise:ch07-35}
证明所有收敛的实数序列空间 $c$ 与所有收敛到零的实数序列空间 $c_{0}$ 关于 $\ell^{\infty}$ 范数是 Banach 空间.
\end{exercise}

\end{document}